\hypertarget{appendix-2ed044358b82}{}\label{referencestart:appendix-2ed044358b82}
\section{Arithmetic policies}\label{appendix-2ed044358b82-arithmetic-policies}

\hypertarget{appendix-2ed044358b82-arithmetic-policies}{}

\textbf{Curriculum depth:} Introduced · \textbf{First seen:} C002 · \textbf{Latest development:} C002

\begin{allmightroledefinition}

\subsection{Definition and mechanism}\label{appendix-2ed044358b82-definition-and-mechanism}

\hypertarget{appendix-2ed044358b82-definition-and-mechanism}{}

Modular, checked, widened, saturated, and arbitrary-precision arithmetic define different outcomes when a mathematical result does not fit.

\end{allmightroledefinition}

A policy states the intended response to an unrepresentable answer: wrap modulo a base, report checked failure, widen storage, clamp to an endpoint, or grow arbitrary precision. Choose by subsystem meaning. Color-byte output and hash arithmetic need different policies.

\subsection{Prerequisites}\label{appendix-2ed044358b82-prerequisites}

\hypertarget{appendix-2ed044358b82-prerequisites}{}

\begin{itemize}
\tightlist
\item
  \hyperref[appendix-c99977d88b33]{Unsigned modular arithmetic} --- Modulo m retains the remainder class after division by m.
\item
  \hyperref[appendix-4fbc03d20829]{Two\textquotesingle s complement} --- For raw unsigned value u in N bits, signed decoding is u when u\textless2ᴺ⁻¹, otherwise u−2ᴺ.
\end{itemize}

\subsection{Used by}\label{appendix-2ed044358b82-used-by}

\hypertarget{appendix-2ed044358b82-used-by}{}

\begin{itemize}
\item
  \hyperref[appendix-3333db47830f]{Fixed-point rescaling} --- uses this mechanism to define its representation, error or validation contract.
\item
  \hyperref[appendix-1e45d0613db3]{Capped weighted updates} --- Capping changes which influence is retained rather than only the representation.
\item
  \href{../07-Vision-XR-and-Playable-Worlds-P02/main.pdf\#nameddest=appendix-fdc7879e8bf2}{Quantization and compressed weights (Vision, XR, and Playable Worlds, Part 2, Reference notes: Quantization and compressed weights)} --- uses arithmetic policies as part of its representation or reasoning.
\item
  \hyperref[appendix-3999c6209649]{Saturating arithmetic} --- uses arithmetic policies as part of its representation or reasoning.
\end{itemize}

\subsection{Recall and next depth}\label{appendix-2ed044358b82-recall-and-next-depth}

\hypertarget{appendix-2ed044358b82-recall-and-next-depth}{}

Explain arithmetic policies without the label, work the example above, then state the assumption whose removal would change the result.

\textbf{Pending:} choosing and documenting policies per subsystem.

\subsection{Source and navigation}\label{appendix-2ed044358b82-source-and-navigation}

\hypertarget{appendix-2ed044358b82-source-and-navigation}{}

\hyperref[01-foundations-i-ch001-arithmetic-becomes-a-policy]{C002 teaching chapter} · \href{../01-Bits-Programs-and-Reliable-Systems-P03/main.pdf\#nameddest=appendix-fd090d982a15}{Part definitions and mechanisms (Bits, Programs, and Reliable Systems, Part 3, Reference notes: Foundations)} · \href{../07-Vision-XR-and-Playable-Worlds-P01/main.pdf\#nameddest=appendix-4ad47bdd269a}{Dependency Map (Vision, XR, and Playable Worlds, Part 1, Reference notes: Dependency Map)}
